\begin{lstlisting}[style=pseudocodigo]
funcion purgarInactivos():
    crear identidadTemp.dat, academicoTemp.dat, contactoTemp.dat

    para cada posicion desde 0 hasta contarRegistros() - 1:
        academico = leerRegistroEnOffset(academico.dat, posicion)

        si academico.marcadoParaPurga != 1:
            identidad = leerRegistroEnOffset(identidad.dat, posicion)
            contacto  = leerRegistroEnOffset(contacto.dat, posicion)
            agregarRegistroAlFinal(identidadTemp.dat, identidad)
            agregarRegistroAlFinal(academicoTemp.dat, academico)
            agregarRegistroAlFinal(contactoTemp.dat, contacto)
        // si marcadoParaPurga es 1, el registro se descarta y no pasa al archivo nuevo
        // (independiente de estadoAcademico: un estudiante en pausa o con permiso
        // no se marca para purga, asi que nunca se elimina solo por su estado academico)

    reemplazarArchivo(identidad.dat, identidadTemp.dat)
    reemplazarArchivo(academico.dat, academicoTemp.dat)
    reemplazarArchivo(contacto.dat, contactoTemp.dat)

    // las posiciones cambiaron: reconstruir los tres mapas desde cero,
    // no se puede reutilizar la lista de recien ingresados anterior
    mapaCarne = reconstruirMapa(identidad.dat, CARNE)
    mapaIdentificacion = reconstruirMapa(identidad.dat, IDENTIFICACION)
    mapaApellido = reconstruirMapa(identidad.dat, APELLIDO)

    retornar EXITO
\end{lstlisting}
